\documentclass[11pt,a4paper]{article}
\usepackage[utf8]{inputenc}
\usepackage[margin=0.75in]{geometry}
\usepackage{hyperref}
\usepackage{fontawesome}
\usepackage{titlesec}
\usepackage{enumitem}

\titleformat{\section}{\Large\bfseries}{\thesection}{1em}{}[\titlerule]
\titlespacing*{\section}{0pt}{*1.5}{*1.2}

% Custom commands
\newcommand{\entry}[4]{
    {\textbf{#1}} \hfill #2 \\
    {\textit{#3}} \hfill \textit{#4} \\
}

\begin{document}

\begin{center}
{\Huge \textbf{Ata Akbari Asanjan}}\\[0.3cm]
\href{mailto:aakbaria@uci.edu}{\faEnvelope\ aakbaria@uci.edu} $\cdot$
\href{https://www.linkedin.com/in/ataakbariasanjan/}{\faLinkedin\ LinkedIn}
\end{center}

\section*{Professional Summary}
% Machine Learning Researcher with 8+ years of experience in developing foundation models, generative architectures, and advanced computer vision systems. Specializes in probabilistic and energy-based deep learning, multi-modal perception, and large-scale training pipelines for Earth and physical sciences. Strong publication record and hands-on expertise in state-of-the-art model design and optimization.
Machine Learning Engineer with 8+ years of experience implementing foundation models, generative architectures, and robust computer vision systems. Specializes in production-grade probabilistic and energy-based deep learning, multi-modal perception pipelines, and scalable ML infrastructure. Proven track record of deploying high-performance models that deliver measurable business impact across diverse domains.
% Machine Learning Engineer with 8+ years of experience in developing and deploying large-scale deep learning systems for computer vision and autonomous systems. Expertise in transformer architectures and multi-modal perception. Proven track record of implementing research and production ML pipelines and leading cross-functional teams to deliver state-of-the-art solutions for complex perception tasks. Strong background in model architecture design, distributed training, and ML system optimization.

\section*{Experience}
\entry{Senior ML Engineer}{Nov. 2023 - Present}{NASA Ames Research Center}{Moffett Field, CA}
\begin{itemize}[nosep,leftmargin=*]
\item Led development of foundation models for Earth atmospheric observation, achieving state-of-the-art performance through transformer architectures and real-time model optimization on terabyte-scale satellite data for complex perception tasks
\item Designed and implemented vector-quantized multi-modal perception models for heliophysics applications, achieving state-of-the-art performance through joint representation learning and real-time optimization across diverse solar observation modalities
\item Implemented probabilistic deep learning architectures for wildfire detection, achieving 10\% improvement over state-of-the-art through uncertainty quantification and ensemble predictions
\item Developed conditional energy-based models with sampling buffers for aviation anomaly detection, delivering 8\% F1-score improvement through robust optimization
\item Led High-Performance Computing (HPC) initiatives to scale machine learning models, implementing multi-GPU training with PyTorch Lightning and DeepSpeed for efficient distributed model training and optimization
% \item Designed and implemented multi-modal transformer architectures with vector quantization for complex perception tasks in heliophysics, enabling joint representation learning across observation modalities
\end{itemize}

\entry{Interdisciplinary ML Engineer}{Nov. 2021 - Nov. 2023}{NASA Ames Research Center}{Moffett Field, CA}
\begin{itemize}[nosep,leftmargin=*]
\item Designed and implemented probabilistic deep learning architectures for wildfire detection, achieving 10\% improvement over state-of-the-art through uncertainty quantification and ensemble predictions
\item Led development of ML-ready datasets for hurricane detection through advanced data fusion techniques, enabling adoption by Federal agencies for operational use
\item Implemented quantum-compatible variational autoencoders using restricted Boltzmann machines with parallel tempering sampling techniques for real-world perception tasks
\item Developed conditional energy-based models with optimized sampling strategies for aviation anomaly detection, delivering 8\% F1-score improvement through robust optimization
\item Led High-Performance Computing initiatives to scale machine learning models across quantum and classical computing infrastructure for efficient distributed training
\end{itemize}

\entry{Research ML Engineer}{Aug. 2019 - Nov. 2021}{NASA Ames Research Center}{Moffett Field, CA}
\begin{itemize}[nosep,leftmargin=*]
\item Designed and delivered introductory and advanced data science curriculum, focusing on Earth science.
\item Developed novel machine learning approaches for anomaly detection in commercial aviation datasets using Interpolation Consistency Training approaches.
\item Developed inverse modeling framework using deep neural networks to infer air pollution concentrations from land surface reflectance measurements, enabling data-poor region monitoring
\item Developed deep learning models for high-resolution coral reef mapping using multi-source satellite imagery through instrument-invariant GAN and Laplacian GAN pyramid super-resolution, achieving high-resolution global coverage
\end{itemize}

\section*{Technical Skills}
\textbf{Programming:} Python (Advanced), SQL, R, Julia, Rust, JavaScript\\
\textbf{ML/AI Frameworks:} PyTorch, TensorFlow, Keras, JAX, Lightning, TIMM, DeepSpeed\\
\textbf{Statistical Methods:} Bayesian Networks, Probabilistic Graphical Models, Uncertainty Quantification, Causal Inference, A/B Testing, Hypothesis Testing, Time Series Analysis, Survival Analysis, Markov Chain Monte Carlo (MCMC), Variational Inference\\
\textbf{Specializations:} Computer Vision, Deep Learning Architecture Design, Foundation Models, Multi-Modal Learning, Probabilistic Deep Learning, Uncertainty Quantification, Time Series Analysis\\
\textbf{Infrastructure:} Google Cloud Platform (GCP), Azure, AWS, Docker, Kubernetes, CI/CD, Terraform, MLflow, Kubeflow, Git, Ray, DVC, Weights \& Biases, Databricks, Spark

\section*{Education}
\entry{Ph.D. in Engineering - Civil}{2019}{University of California, Irvine}{}\entry{M.Sc. in Engineering - Civil}{2016}{University of California, Irvine}{}\entry{B.Sc. in Engineering - Civil}{2014}{University of Tabriz}{}

\vspace{-2ex}
\section*{Selected Publications}
\begin{itemize}[nosep,leftmargin=*]
\item Esencan, M., Kumar, T. A., \textbf{Akbari Asanjan, A.}, Lott, P. A., Mohseni, M., Unlu, C., ... \& Ho, A. (2024). ``Combinatorial reasoning: selecting reasons in generative AI pipelines via combinatorial optimization.'' \textit{arXiv preprint arXiv:2407.00071}.
\item \textbf{Akbari Asanjan, A.}, Memarzadeh, M., Lott, P. A., Rieffel, E., \& Grabbe, S. (2023). ``Probabilistic wildfire segmentation using supervised deep generative model from satellite imagery.'' \textit{Remote Sensing}, 15(11), 2718.
\item Szwarcman, D., Roy, S., Fraccaro, P., Gíslason, Þ. E., Blumenstiel, B., Ghosal, R., \textbf{Akbari Asanjan, A.}, ... \& Moreno, J. B. (2024). ``Prithvi-EO-2.0: A versatile multi-temporal foundation model for earth observation applications.'' \textit{arXiv preprint arXiv:2412.02732}.
\item \textbf{Akbari Asanjan, A.}, Das, K., Li, A., Chirayath, V., Torres-Perez, J., \& Sorooshian, S. (2020). ``Learning instrument invariant characteristics for generating high-resolution global coral reef maps.'' In \textit{Proceedings of the 26th ACM SIGKDD International Conference on Knowledge Discovery \& Data Mining} (pp. 2617-2624).
% \item \textbf{Akbari Asanjan, A.}, Brady, L., Suri, N., Izquierdo, Z. G., Lott, P. A., Memarzadeh, M., ... \& Grabbe, S. (2023, June). ``\textit{Image-to-Image Wildfire Detection via Quantum-Compatible Variational Segmentation from Remotely-sensed Data}.'' In \textit{Earth Science Technology Forum}.
\end{itemize}

\section*{Selected Projects}
\begin{itemize}[nosep,leftmargin=*]

\item \textbf{Generative AI \& Computer Vision:}
  \begin{itemize}[nosep]
    \item Developed text-guided facial editing system using latent Diffusion Transformer (DiT) architecture, enabling precise attribute manipulation through natural language
    \item Implemented custom CLIP-based text-image matching system for automated training data curation on CelebA-HQ
    \item Developed efficient training pipeline using PyTorch Lightning for distributed training
    \item Created interactive demo using Gradio for real-time facial attribute editing
  \end{itemize}
    
% \item \textbf{Web Development \& Geospatial:}
%   \begin{itemize}[nosep]
%     \item Developed Google-Earth-like satellite imagery visualization platform using Flask tile-server, React, and Three.js for interactive 3D rendering
%     \item Implemented WebGL-based 3D terrain rendering with custom shaders and texture mapping for high-performance visualization
%     \item Built RESTful API services using FastAPI and Redis caching for efficient tile delivery and data streaming
%     \item Created scalable data pipeline using GeoServer, PostGIS, and GDAL for processing and serving satellite imagery fast and scalable.
%   \end{itemize}

\item \textbf{Computer Vision \& Robotics:}
  \begin{itemize}[nosep]
    \item Developed UAV GPS coordinate localization system using NeRF-based scene reconstruction and matching
    \item Implemented comparative analysis with traditional SfM and SLAM approaches using OpenCV and COLMAP
    \item Built real-time pose estimation pipeline using PyTorch3D and custom differentiable rendering
    \item Created efficient point cloud registration using CUDA-accelerated ICP and feature matching
  \end{itemize}

% \item \textbf{Natural Language Processing \& Data Visualization:}
%   \begin{itemize}[nosep]
%     \item Developed RAG-based weather analysis system using Llama model and structured weather data
%     \item Created React-based dashboard integrating weather summaries with interactive geospatial visualizations
%     \item Implemented efficient data pipeline for real-time weather data processing and visualization
%   \end{itemize}

\end{itemize}

\end{document} 